\documentclass[12pt,a4paper]{amsart}

\usepackage{amsmath,amssymb,amsthm}
\usepackage{mathtools}
\usepackage{enumitem}
\usepackage{hyperref}
\usepackage{booktabs}
\usepackage{tikz-cd}
\usetikzlibrary{positioning,arrows.meta}

%%%─ Environments
\newtheorem{theorem}{Theorem}[section]
\newtheorem{lemma}[theorem]{Lemma}
\newtheorem{proposition}[theorem]{Proposition}
\newtheorem{corollary}[theorem]{Corollary}
\theoremstyle{definition}
\newtheorem{definition}[theorem]{Definition}
\newtheorem{hypothesis}[theorem]{Hypothesis}
\newtheorem{remark}[theorem]{Remark}

%%%─ Macros
\newcommand{\cA}{\mathcal{A}}
\newcommand{\cAg}{\mathcal{A}_\gamma}
\newcommand{\R}{\mathbb{R}}
\newcommand{\N}{\mathbb{N}}
\newcommand{\cK}{K_{B_c}}
\newcommand{\cKA}{K_{\cA}}
\newcommand{\normh}[2]{\|#1\|_{#2}}
\newcommand{\SK}{S_K}
\newcommand{\AK}{A_K}
\newcommand{\AKlin}{A_K^{\mathrm{lin}}}
\newcommand{\AKquad}{A_K^{\mathrm{quad}}}
\newcommand{\BK}{B_K}
\newcommand{\Kopt}{K^*}
\newcommand{\Clin}{C_{\mathrm{lin}}(M,K)}
\newcommand{\Cquad}{C_{\mathrm{quad}}(M,K)}
\newcommand{\Ctwo}{C_2(\alpha)}
\newcommand{\Cinfty}{C_\infty}

\subjclass[2020]{34B24, 47B34, 34E20, 47G10, 34L20}
\keywords{Sturm--Liouville operators, turning-point asymptotics,
  spectral growth exponent, two-scale truncation, logarithmic rate,
  spectral-geometric invariant, universality, layer-survival,
  spectrally regularized WKB class, variational truncation,
  Langer error, discrete floor projection, iterated asymptotics,
  diagonal substitution, stability of optimal truncation,
  substitution stability principle}

\begin{document}

\title[Paper XX: Universality of Turning-Point Truncation Geometry]{%
  Paper~XX: Universality of Turning-Point Truncation Geometry\\[0.3em]
  \large Spectral-Geometric Invariants and Two-Scale Logarithmic Losses}

\author{\textit{prolate-gram-coercivity programme}}
\date{May 2026 --- FROZEN}

\begin{abstract}
For operators $\cAg\in\mathfrak{T}_\gamma$
(spectrally regularized WKB class, Definition~\ref{def:opclass}),
the Langer term $\AK = \AKlin + \AKquad$ satisfies,
under the pointwise Langer scaling law
$\|e_j\|_{C^0}\leq C_0 j^\gamma c^{-\beta}$
(Hypothesis~\ref{hyp:langer_pointwise}),
\[
  \|\AKlin\| \leq C_{\rm lin}\,K^{1+\gamma}c^{-\beta},
  \qquad
  \|\AKquad\| \leq C_{\rm quad}\,K^{1+2\gamma}c^{-2\beta}
\]
for all $K\geq K_0$, $c\geq c_0$, with constants $C_{\rm lin}, C_{\rm quad}$
independent of $K$ and $c$
(Lemma~\ref{lem:trunc_universal}).
The optimal truncation $\Kopt$ is the unique minimizer of the
fixed-$c$ functional $F(K;c) := L\,K^{1+\gamma}c^{-\beta} + T\,e^{-\alpha K/2}$
over the window $K\in[c_1\log c,c_2\log c]$
(Lemma~\ref{lem:kopt_variational}), satisfying
$K^*(c) = \frac{2\beta}{\alpha}\log c + \frac{2\gamma}{\alpha}\log\log c + O(1)$.
The transition from fixed-parameter bounds to the diagonal
$K = K^*(c)$ is licensed by the Substitution Stability Principle
(Proposition~\ref{prop:substitution_stability}),
which is the unique formal bridge between the fixed-parameter and
diagonal regimes; logarithmic instability under exponentiation
is controlled by an explicit slowly-varying substitution estimate,
uniform in $\gamma$ on compact subsets of $(0,\infty)$.
All implied constants are $c$-admissible with no hidden dependence
(Remark~\ref{rem:cadmissible}, Lemma~\ref{lem:cadm_closure}).
All results are conditional on H1--H2.
\end{abstract}

\maketitle
\tableofcontents
\bigskip\hrule\bigskip

%% ============================================================
\section{Introduction and Logical Status}
\label{sec:intro}
%% ============================================================

\begin{center}
\renewcommand{\arraystretch}{1.3}
\begin{tabular}{l l l}
\toprule
\textbf{Paper} & \textbf{Content} & \textbf{Logical status} \\
\midrule
XVIII & BR3: qualitative Airy universality
  & Theorem (Olver + Slepian) \\
XIX   & Two-scale rate; lin/quad split
  & Theorem (Airy; H1--H2 verified) \\
XX    & Universal rate; Layer-Survival; variational $\Kopt$
  & Conditional theorem ($\mathfrak{T}_\gamma$, H1--H2) \\
\bottomrule
\end{tabular}
\end{center}

\begin{remark}[Scope of XX]
\label{rem:logstatus}
All results are theorems within $\mathfrak{T}_\gamma$.
H1--H2 are imposed axioms, not derived for general $V$.
Realization: verified for Airy; sketch for power family; open in general.
\end{remark}

\begin{remark}[$c$-admissible constants]
\label{rem:cadmissible}
We call a constant \emph{$c$-admissible} if it is independent of
the approximation parameter $c$.
There are four sources of $c$-admissibility in this paper:
\begin{itemize}[nosep]
  \item \textbf{Axiomatic:} $C_0$, $C_\infty$ declared $c$-free
    by H1 and H2 (operator-intrinsic).
  \item \textbf{Analytic:} $C_\sigma = \frac{|B_2|}{12}\sigma(\sigma-1)$
    depends only on $\sigma$, not on $K$ or $c$
    (eq.~\eqref{eq:EM_remainder}).
  \item \textbf{Derived:} products and finite sums of $c$-admissible
    constants are $c$-admissible (Lemma~\ref{lem:cadm_closure});
    covers $C_{\rm lin}$, $C_{\rm quad}$, $C_{\sigma,1}$,
    $C_{\sigma,2}$, and all layer constants $C_n$.
  \item \textbf{Operator-intrinsic:}
    $C_2(\alpha)$ is $c$-admissible by definition
    (Remark~\ref{rem:BK_definition}).
\end{itemize}
Every constant appearing in this paper falls under one of these
four cases.
No intermediate estimate introduces hidden $c$-dependence;
all constants are traceable to one of the four declared sources above.
The chain $C_0,C_\infty \to C_{\rm lin},C_{\rm quad} \to C_{\rm tail}$
is therefore completely $c$-free.
\end{remark}

\begin{lemma}[$c$-admissible closure]
\label{lem:cadm_closure}
The class of $c$-admissible constants is closed under
finite products and finite sums.
\end{lemma}

\begin{proof}
If $A$ and $B$ are independent of $c$, so are $AB$ and $A+B$.
\qed
\end{proof}

\begin{remark}[Two sources of $c$-dependence, strictly separated]
\label{rem:dag_csplit}
In this paper, $c$ appears in two logically distinct roles:
\begin{itemize}[nosep]
  \item \textbf{Operator estimates (fixed-parameter regime):}
    $c$ appears only as a scaling factor $c^{-\beta}$;
    all constants multiplying these powers are $c$-admissible;
    $K$ and $c$ are independent in this regime.
  \item \textbf{Domain restriction (diagonal regime):}
    $c$ appears in the sandwich window $[c_1\log c, c_2\log c]$
    and in $K^*(c)$ via the fixed-$c$ FOC.
    This is $c$-dependence in the \emph{domain of optimization},
    not in the operator geometry.
\end{itemize}
These two roles are disjoint.
Domain dependence on $c$ does not introduce $c$-dependence
into any operator constant.
The formal bridge between the two roles is
Proposition~\ref{prop:substitution_stability}.
\end{remark}

\begin{remark}[Two-regime structure and quantifier discipline]
\label{rem:diagonal_regime}
The paper is organized in two strictly separated regimes
(see Remark~\ref{rem:dag_csplit} for the precise separation):
\begin{itemize}[nosep]
  \item \textbf{Fixed-parameter regime.}
    All operator inequalities hold for all $K\geq K_0$, $c\geq c_0$
    independently; all implied constants are $c$-admissible.
  \item \textbf{Diagonal regime.}
    $\forall c\geq c_0,\;\exists!\,K^*(c)\in[c_1\log c, c_2\log c]$;
    asymptotics derived as $c\to\infty$.
\end{itemize}
The transition is performed by exactly one proposition
(Proposition~\ref{prop:substitution_stability}),
which licenses all diagonal evaluations.
\end{remark}

%% ============================================================
\section{Substitution Stability Principle}
\label{sec:substprin}
%% ============================================================

\begin{proposition}[Substitution Stability Principle]
\label{prop:substitution_stability}
Let
\[
  |E(K,c)| \leq \sum_{n=1}^N C_n\,K^{a_n}c^{-b_n}
\]
hold for all $K \geq K_0$, $c \geq c_0$, where all $C_n$
are $c$-admissible (Remark~\ref{rem:cadmissible}).

Let $K(c)$ be any function satisfying
\begin{equation}\label{eq:Kc_window}
  K(c) \in [c_1 \log c,\, c_2 \log c]
  \quad \text{for all sufficiently large } c.
\end{equation}

Then:
\begin{itemize}[nosep]
  \item[(i)] The bound remains valid under substitution
    $K \mapsto K(c)$:
    \[
      |E(K(c),c)| \leq \sum_{n=1}^N C_n\,(K(c))^{a_n}c^{-b_n}.
    \]
    This is immediate since the bound holds pointwise in $K$
    for all $K \geq K_0$; no uniformity in $K$ beyond
    pointwise validity is required.
  \item[(ii)] Each term $(K(c))^{a_n}c^{-b_n}$ admits an explicit
    asymptotic order via the window control~\eqref{eq:Kc_window}:
    \[
      (K(c))^{a_n} = O(\log^{a_n} c),
    \]
    and for $K(c) \sim \tfrac{2\beta}{\alpha}\log c$ the precise
    slowly-varying correction is given by~\eqref{eq:Kstar_power}.
  \item[(iii)] Asymptotic dominance relations among terms
    are those of Propositions~\ref{prop:asymp_order}
    and~\ref{prop:scale_algebra}, evaluated termwise.
\end{itemize}

In particular, every asymptotic comparison in the diagonal regime
is a reduction of a fixed-parameter bound under this principle;
no separate uniformity argument is needed.
\end{proposition}

\begin{proof}
(i) is pointwise substitution into a pointwise-valid inequality.
(ii) follows from $K(c) \leq c_2\log c$ (upper sandwich bound).
(iii) is Propositions~\ref{prop:asymp_order}--\ref{prop:scale_algebra}.
\qed
\end{proof}

%% ============================================================
\section{Asymptotic Scale and Regime-Dependent Order}
\label{sec:asymp}
%% ============================================================

\begin{definition}[Asymptotic scale representatives]
\label{def:scale}
Let $\mathcal{S} = \{c^{-s}(\log c)^t : s>0,\, t\in\R\}$,
used as a representative system of asymptotic equivalence
classes: $f\sim g \Leftrightarrow f/g\to 1$ as $c\to\infty$.
The total order $\succ$ on equivalence classes:
$c^{-s_1}(\log c)^{t_1}\succ c^{-s_2}(\log c)^{t_2}$
iff $s_1<s_2$, or ($s_1=s_2$ and $t_1>t_2$).
When we write $f = O(g)$ or $f\prec g$
for $f,g\in\mathcal{S}$, this is a statement about the
asymptotic ratio $f/g$ as $c\to\infty$;
no algebraic structure on $\mathcal{S}$ is assumed.
\end{definition}

\begin{remark}[Dominance is regime-dependent, not algebraic]
\label{rem:scale_regime}
All asymptotic comparisons ($\prec$, $O(\cdot)$, $\sim$)
in this paper are evaluated in the diagonal regime
$K = K^*(c) \sim \frac{2\beta}{\alpha}\log c$.
Dominance is verified explicitly for each pair;
not assumed as a global property of $\mathcal{S}$.
\end{remark}

\begin{proposition}[Multiplication on $\mathcal{S}$]
\label{prop:scale_algebra}
For $f = c^{-s_1}(\log c)^{t_1}$,
$g = c^{-s_2}(\log c)^{t_2}\in\mathcal{S}$:
$f\cdot g = c^{-(s_1+s_2)}(\log c)^{t_1+t_2}\in\mathcal{S}$.
If $f/g\to\infty$: $f+g = O(f)$ with constant $2$.
\end{proposition}

\begin{proposition}[Layer dominance in diagonal regime]
\label{prop:asymp_order}
With $K = K^*(c)\sim\frac{2\beta}{\alpha}\log c$, for $n\geq 2$:
$c^{-n\beta}(\log c)^{1+n\gamma} \prec c^{-\beta}(\log c)^{1+\gamma}$.
\end{proposition}

%% ============================================================
\section{Logical Dependency DAG}
\label{sec:dag_logical}
%% ============================================================

\begin{remark}[DAG structure]
\label{rem:dag_reading}
The following diagram is an organizational dependency map
of the proof structure.
It is not a mathematical object or theorem;
acyclicity is a description of proof discipline,
not a separately verified combinatorial claim.
Solid arrows: direct logical implication used in the proofs.
All couplings between $K$ and $c$ pass exclusively through
Proposition~\ref{prop:substitution_stability}
and Lemma~\ref{lem:diagonal_sub}.
\end{remark}

\[
\begin{tikzcd}[row sep=2.2em, column sep=1.8em,
  every arrow/.append style={-Latex}]
\boxed{\mathrm{H1}} \ar[dr]
& & \boxed{\mathrm{H2}} \ar[dl]
& & \boxed{F(K;c)\text{ fixed-}c} \ar[d] \\
& \text{lem:trunc\_univ} \ar[d]
& & & \text{lem:kstar\_to\_infty} \ar[d] \\
& \text{thm:rate} \ar[dr]
& & & \text{lem:loglog} \ar[d] \ar[dr] \\
& & \text{\textbf{prop:subst}} \ar[d]
& \text{lem:kopt(ii)} \ar[l]
& \text{lem:stability} \ar[l]
& \text{(ii') reg.} \ar[l] \\
& & \text{lem:diag\_sub} \ar[d] & & & \\
& & \text{cor:rate} & & &
\end{tikzcd}
\]

\begin{remark}[Acyclicity guarantees]
\label{rem:dag_nocycle}
\begin{itemize}[nosep]
  \item $F(K;c)$ defined with $K,c$ independent;
    $K^*(c)$ extracted by $F'(\,\cdot\,;c)=0$ at fixed $c$:
    one-directional, no feedback.
  \item Lemma~\ref{lem:loglog} uses only $K^*\to\infty$
    and sign of $F'$; does not use~\eqref{eq:kopt_iterated}.
  \item Proposition~\ref{prop:substitution_stability}
    uses only the pointwise-valid fixed-parameter bound
    and the window control~\eqref{eq:Kc_window};
    it is the unique formal bridge between regimes.
  \item The domain window $[c_1\log c, c_2\log c]$ introduces
    $c$-dependence only in the optimization domain,
    not in any operator constant
    (Remark~\ref{rem:dag_csplit}).
  \item (ii') certifies self-consistency only;
    produces no new asymptotic information.
\end{itemize}
\end{remark}

%% ============================================================
\section{Framework}
\label{sec:framework}
%% ============================================================

\begin{definition}[$\mathfrak{T}_\gamma$]
\label{def:opclass}
$\mathfrak{T}_\gamma$ ($\gamma,\beta>0$): self-adjoint
$\cAg=-\partial_\xi^2+V(\xi)$ on $L^2(\R_+)$,
$V\geq 0$, $V\to+\infty$, discrete spectrum,
satisfying (i) $a_j\sim C_\gamma j^\gamma$;
(ii) $V\in C^2$, simple turning points; (iii) H1; (iv) H2.
The parameter $c$ is an approximation scale parameter,
not a spectral parameter of $\cAg$;
operators in $\mathfrak{T}_\gamma$ are defined independently of $c$.
All operator-intrinsic objects (eigenfunctions, spectral projections,
spectral remainder norms) are $c$-free at the level of definition.
\end{definition}

\begin{hypothesis}[H1: Pointwise Langer scaling]
\label{hyp:langer_pointwise}
$\exists\,C_0<\infty$ (independent of $j,K,c$):
$\|e_j\|_{C^0([0,M])}\leq C_0\,j^\gamma c^{-\beta}$
for $c\geq c_0(j)$.
\emph{Status:} Airy: XVIII;
power family: \texttt{olver\_growth\_lemma.tex};
general: H1-K.
\end{hypothesis}

\begin{hypothesis}[H2: $C_\infty$ boundedness]
\label{hyp:Cinfty}
$C_\infty=\sup_{j\geq 1}\|u_j\|_{C^0([0,M])}<\infty$,
where $u_j$ are eigenfunctions of $\cAg$
(operator-intrinsic, independent of $c$).
\emph{Status:} Airy: standard; power family: WKB; general: open.
\end{hypothesis}

%% ============================================================
\section{Operator Decomposition}
\label{sec:opdecomp}
%% ============================================================

\begin{remark}[Definition-level $c$-independence of $\BK$]
\label{rem:BK_definition}
$\BK$ is the spectral remainder of $\cAg$ after truncation
at index $K$: a functional of the spectral projections
$\{P_j : j > K\}$ of $\cAg$ alone.
No approximation parameter $c$ enters the definition of $\BK$.
Consequently:
\begin{itemize}[nosep]
  \item $C_2(\alpha) = \sup_{K\geq 0}\|\BK\|^2 e^{\alpha K}$
    is a constant of $\cAg$ and $\alpha$;
    it is $c$-admissible by definition, not merely by argument.
  \item The supremum over all $K$ is compatible with evaluation
    at any specific $K(c)$: the bound $\|\BK\|\leq C_2^{1/2}e^{-\alpha K/2}$
    holds pointwise in $K$, and $K^*(c)$ is simply one choice of $K$.
    No uniformity in $K$ beyond pointwise validity is required.
\end{itemize}
\end{remark}

\begin{lemma}[$\gamma$-Agnostic nodes]
\label{lem:gamma_agnostic}
Without H1--H2, for all $\cAg\in\mathfrak{T}_\gamma$:
\begin{itemize}[nosep]
\item[(i)] $\SK-\cKA=\AK-\BK$;
\item[(ii)] tensor norm bound;
\item[(iii)] Bessel tail;
\item[(iv)] Tonelli/DCT;
\item[(v)] $\|\BK\|\leq\Ctwo^{1/2}e^{-\alpha K/2}$,
  $C_2(\alpha) := \sup_{K\geq 0}\|\BK\|^2 e^{\alpha K}$;
  $c$-admissible by Remark~\ref{rem:BK_definition}.
\item[(vi)] Kato norm-resolvent.
\end{itemize}
\end{lemma}

%% ============================================================
\section{Truncation Geometry}
\label{sec:trunc}
%% ============================================================

\begin{lemma}[Universal truncation geometry, fixed-parameter]
\label{lem:trunc_universal}
Under H1--H2, with
$C_{\rm lin} = \frac{2MC_\infty C_0}{1+\gamma}$,
$C_{\rm quad} = \frac{MC_0^2}{1+2\gamma}$
(both $c$-admissible by Lemma~\ref{lem:cadm_closure}),
for all $K\geq 1$, $c\geq c_0$:
\begin{align}
  \normh{\AKlin}{}
    &\leq C_{\rm lin}\,K^{1+\gamma}c^{-\beta}
    + C_{\sigma,1}\,K^\gamma c^{-\beta},\label{eq:AKlin}\\
  \normh{\AKquad}{}
    &\leq C_{\rm quad}\,K^{1+2\gamma}c^{-2\beta}
    + C_{\sigma,2}\,K^{2\gamma}c^{-2\beta},\label{eq:AKquad}
\end{align}
$C_{\sigma,1} = \tfrac{1}{2}C_0 C_\infty$,
$C_{\sigma,2} = \tfrac{1}{2}C_0^2$
($c$-admissible, Lemma~\ref{lem:cadm_closure}).
\end{lemma}

\begin{proof}
Apply Euler--Maclaurin to $\sum_{j=1}^K j^\sigma$:
\begin{equation}\label{eq:EM}
  \sum_{j=1}^K j^\sigma
  = \frac{K^{1+\sigma}}{1+\sigma} + \frac{1}{2}K^\sigma + R_\sigma(K),
\end{equation}
\begin{equation}\label{eq:EM_remainder}
  |R_\sigma(K)| \leq C_\sigma K^{\sigma-1},
  \quad
  C_\sigma = \frac{|B_2|}{12}\sigma(\sigma-1),
  \quad B_2=\tfrac{1}{6}.
\end{equation}
$C_\sigma$ depends only on $\sigma$, not on $K$ or $c$.
Multiply by $c$-admissible prefactors from H1--H2;
closure by Lemma~\ref{lem:cadm_closure}.
\end{proof}

%% ============================================================
\section{Variational Optimal Truncation}
\label{sec:kopt}
%% ============================================================

\begin{lemma}[$K^*(c)\to\infty$]
\label{lem:kstar_to_infty}
For each fixed $c$, $F(\,\cdot\,;c)=LK^{1+\gamma}c^{-\beta}+Te^{-\alpha K/2}$
has a unique minimizer $K^*(c)$; $K^*(c)\to\infty$ as $c\to\infty$.
\end{lemma}

\begin{proof}
$F'(K_0;c)<0$ for large $c$ at any fixed $K_0$.
\end{proof}

\begin{lemma}[Sandwich on $K^*(c)$]
\label{lem:loglog}
With $c_1=2\beta/\alpha$, $c_2=2(\beta+1)/\alpha$,
for all sufficiently large $c$:
\begin{equation}\label{eq:sandwich}
  c_1\log c\leq K^*(c)\leq c_2\log c,
  \qquad \log K^*(c)=\log\log c+O(1).
\end{equation}
\end{lemma}

\begin{lemma}[Variational $K^*$]
\label{lem:kopt_variational}

\textup{(i)} \textbf{Fixed-$c$ existence and uniqueness.}
For each fixed $c$, $F''(K;c)>0$ on $[c_1\log c, c_2\log c]$
(strict convexity; $F''\to 0$ pointwise, no uniform-in-$c$ claim);
unique minimizer $K^*(c)$ satisfies
\begin{equation}\label{eq:FOC}
  L(1+\gamma){K^*}^\gamma c^{-\beta}=\tfrac{\alpha}{2}T e^{-\alpha K^*/2}.
\end{equation}

\textup{(ii)} \textbf{Iterated asymptotics.}
\begin{equation}\label{eq:log_FOC}
  \tfrac{\alpha}{2}K^*=\beta\log c+\gamma\log K^*+C_1.
\end{equation}
Three-step derivation (no circularity):
Step~1: $K^*\to\infty$ gives leading term.
Step~2: sandwich gives $\log K^*=\log\log c+O(1)$.
Step~3:
\begin{equation}\label{eq:kopt_iterated}
  K^*(c)=\frac{2\beta}{\alpha}\log c+\frac{2\gamma}{\alpha}\log\log c+O(1).
\end{equation}

\textup{(ii')} \textbf{Role of $\lambda(c)$ (self-consistency only).}
The iteration map $\Phi_c(K)=\frac{2\beta}{\alpha}\log c+\frac{2\gamma}{\alpha}\log K+C_1$
has Lipschitz constant $\lambda(c) = \frac{2\gamma}{\alpha K^*(c)}\to 0$,
hence $\lambda(c)<\tfrac{1}{2}$ for large $c$.
\emph{Role:} self-consistency certificate only.
Existence and uniqueness of $K^*(c)$ follows from (i),
not from the Banach fixed-point theorem.
No conclusion depends on $\lambda(c)$ being small uniformly in $c$.

\textup{(iii)} $F(K^*;c)=O(c^{-\beta}(\log c)^{1+\gamma})$.
\end{lemma}

\begin{lemma}[Stability of asymptotics]
\label{lem:stability_kopt}
An $O(1)$ perturbation of~\eqref{eq:log_FOC} shifts $K^*$ by $O(1)$.
\end{lemma}

%% ============================================================
\section{Diagonal Substitution}
\label{sec:diagonal}
%% ============================================================

\begin{lemma}[Diagonal substitution: explicit evaluation]
\label{lem:diagonal_sub}
This lemma applies Proposition~\ref{prop:substitution_stability}
to the specific choice $K(c) = K^*(c)$.

\emph{Input bound} (Theorem~\ref{thm:universal_rate},
all constants $c$-admissible):
\begin{equation}\label{eq:fixed_param_bound}
  \normh{\cK-\cKA}{}
  \leq C_{\rm lin}\,K^{1+\gamma}c^{-\beta}
  + C_{\rm quad}\,K^{1+2\gamma}c^{-2\beta}
  + C_{\rm tail}\,e^{-\alpha K/2}.
\end{equation}

\emph{Slowly-varying substitution stability
(uniform in $\gamma$ on compact sets).}
Let $h(c) := K^*(c)\big/\bigl(\tfrac{2\beta}{\alpha}\log c\bigr)$.
By~\eqref{eq:kopt_iterated}: $h(c)\to 1$,
$|h(c)-1|=O(\log\log c/\log c)$.
By MVT $(1+u)^{1+\gamma}=1+(1+\gamma)\theta^\gamma u$
with $|u|\leq\tfrac{1}{2}$ for large $c$,
the $O$-constant in
\begin{equation}\label{eq:hgamma}
  h(c)^{1+\gamma} = 1 + O\!\left(\frac{\log\log c}{\log c}\right)
\end{equation}
is bounded by $(1+\gamma)2^\gamma$,
finite on every compact $[\gamma_0,\gamma_1]\subset(0,\infty)$.
No uniformity as $\gamma\to\infty$ is claimed.
Hence
\begin{equation}\label{eq:Kstar_power}
  (K^*(c))^{1+\gamma}
  = \Bigl(\tfrac{2\beta}{\alpha}\log c\Bigr)^{1+\gamma}
  \Bigl(1 + O\Bigl(\tfrac{\log\log c}{\log c}\Bigr)\Bigr),
\end{equation}
uniformly in $\gamma$ on compact subsets of $(0,\infty)$.

\emph{Conclusion}
(by Proposition~\ref{prop:substitution_stability}(i)--(iii)):
\begin{itemize}[nosep]
  \item[(i)] $K^*(c)^{1+\gamma}c^{-\beta}=O(c^{-\beta}(\log c)^{1+\gamma})$.
  \item[(ii)] $K^*(c)^{1+2\gamma}c^{-2\beta}
    \prec c^{-\beta}(\log c)^{1+\gamma}$.
  \item[(iii)] $C_{\rm tail}\,e^{-\alpha K^*(c)/2}
    \prec c^{-\beta}(\log c)^{1+\gamma}$.
  \item[(iv)] $\normh{\cK-\cKA}{}
    =O(c^{-\beta}(\log c)^{1+\gamma})$;
    linear term dominant.
\end{itemize}
\end{lemma}

\begin{remark}[Discrete projection]
\label{rem:kopt_discrete}
$K_{\rm int}=\lfloor K^*\rceil$:
$F(K_{\rm int})-F(K^*)\prec c^{-\beta}(\log c)^{1+\gamma}$.
\end{remark}

%% ============================================================
\section{Layer-Survival Selection Principle}
\label{sec:layersurvival}
%% ============================================================

\begin{theorem}[Layer-Survival]
\label{thm:layer_survival}
For $n\geq 1$, under H1, all $K\geq K_0$, $c\geq c_0$:
$\normh{A_K^{(n)}}{}\leq C_n K^{1+n\gamma}c^{-n\beta}$
with $c$-admissible $C_n$ (Lemma~\ref{lem:cadm_closure}).
Diagonal evaluation by
Proposition~\ref{prop:substitution_stability}:
\[
  \normh{A_K^{(n)}}{}_{K=K^*(c)}
  = O(c^{-n\beta}(\log c)^{1+n\gamma})
  \prec c^{-\beta}(\log c)^{1+\gamma}\quad(n\geq 2).
\]
\end{theorem}

%% ============================================================
\section{Main Theorems}
\label{sec:main}
%% ============================================================

\begin{theorem}[Two-scale rate, fixed-parameter]
\label{thm:universal_rate}
For $\cAg\in\mathfrak{T}_\gamma$, all $K\geq K_0$, $c\geq c_0$:
\[
  \normh{\cK-\cKA}{}
  \leq C_{\rm lin}\,K^{1+\gamma}c^{-\beta}
  + C_{\rm quad}\,K^{1+2\gamma}c^{-2\beta}
  + \Ctwo^{1/2}e^{-\alpha K/2},
\]
all constants $c$-admissible
(Remark~\ref{rem:cadmissible}, Lemma~\ref{lem:cadm_closure}).
\end{theorem}

\begin{corollary}[Universal logarithmic rate]
\label{cor:universal_lograte}
Apply Proposition~\ref{prop:substitution_stability}
and Lemma~\ref{lem:diagonal_sub} at $K=\lfloor K^*(c)\rceil$:
\[
  \normh{\cK-\cKA}{}=O(c^{-\beta}(\log c)^{1+\gamma}).
\]
Leading constant $C_{\rm lin}(2\beta/\alpha)^{1+\gamma}$
($c$-admissible).
\end{corollary}

\begin{remark}[Spectral-geometric invariants]
\label{rem:invariants}
$1+\gamma$: leading exponent; $1+2\gamma$: subleading.
$\gamma$ governs the log-exponent and log-log correction
but is absent from the leading coefficient $2\beta/\alpha$ of $K^*$.
\end{remark}

\begin{remark}[O3: sharpness and axiom status]
\label{rem:O3_sharpness}
O3a (lower bound, independent of H1--H2): open.
O3b (uniform lower bound, would need H3): open; H3 not imposed.
Neither blocks any upper-bound conclusion.
\end{remark}

%% ============================================================
\section{Model Family}
\label{sec:models}
%% ============================================================

\begin{center}
\renewcommand{\arraystretch}{1.4}
\begin{tabular}{l l c c c c c}
\toprule
\textbf{Model} & $V$ & $\gamma$ & $\beta$
  & $1{+}\gamma$ & $K^*$ leading & H1 \\
\midrule
Airy     & $\xi$   & $2/3$ & $1/3$ & $5/3$
  & $\frac{2}{3\alpha}\log c$ & Proved \\
Harmonic & $\xi^2$ & $1$   & $1/2$ & $2$
  & $\frac{1}{\alpha}\log c$ & Sketch \\
$|\xi|^\nu$ & & $\frac{2\nu}{\nu+2}$ & $\frac{\nu}{\nu+2}$
  & $\frac{3\nu+2}{\nu+2}$
  & $\frac{2\nu}{(\nu+2)\alpha}\log c$ & Sketch \\
General $V$ & & $\gamma$ & $\beta$ & $1+\gamma$
  & $\frac{2\beta}{\alpha}\log c$ & via H1-K \\
\bottomrule
\end{tabular}
\end{center}

%% ============================================================
\section{Open Tasks}
\label{sec:open}
%% ============================================================

\begin{center}
\renewcommand{\arraystretch}{1.3}
\begin{tabular}{l l l}
\toprule
\textbf{ID} & \textbf{Statement} & \textbf{Status} \\
\midrule
O1 & H1 general $V$ via H1-K & closed (conditional) \\
O2 & H2 for $\nu\neq 1$ & open \\
O3a & lower bound (indep.\ H1--H2) & open \\
O3b & uniform lower bound (would need H3) & open \\
O4 & Olver-Watson sketch$\to$proof & open \\
\hline
\multicolumn{2}{l}{O2, O3, O4 non-blocking} & \checkmark \\
\bottomrule
\end{tabular}
\end{center}

%% ============================================================
\section{Programme Conclusion}
\label{sec:conclusion}
%% ============================================================

\textbf{XVIII:} Qualitative BR3.\quad
\textbf{XIX:} Airy rate.\quad
\textbf{XX:} Universal rate $O(c^{-\beta}(\log c)^{1+\gamma})$:
\begin{itemize}[nosep,leftmargin=1.5em]
  \item $c$-admissible constant chain, four sources,
    no hidden dependence sealed
    (Remark~\ref{rem:cadmissible}, Lemma~\ref{lem:cadm_closure});
  \item $B_K$ and $C_2(\alpha)$ $c$-free by definition
    (Remark~\ref{rem:BK_definition});
  \item domain vs.\ operator $c$-dependence strictly separated
    (Remark~\ref{rem:dag_csplit});
  \item Substitution Stability Principle as unique formal bridge
    between fixed-parameter and diagonal regimes
    (Proposition~\ref{prop:substitution_stability});
  \item fixed-parameter bounds uniform via explicit EM remainder
    (eq.~\eqref{eq:EM_remainder});
  \item slowly-varying stability uniform in $\gamma$ on compact
    subsets; $O$-constant $(1+\gamma)2^\gamma$ explicit;
    no claim for $\gamma\to\infty$;
  \item $\lambda(c)<\tfrac{1}{2}$: self-consistency certificate,
    not existence argument
    (Lemma~\ref{lem:kopt_variational}(ii'));
  \item DAG is organizational map, not mathematical theorem
    (Remark~\ref{rem:dag_reading});
  \item Layer-Survival; invariants $1+\gamma$, $1+2\gamma$.
\end{itemize}

\bigskip
\begin{quote}\itshape\large
The leading logarithmic exponent is selected not by maximal
$K$-growth, but by the highest interaction layer surviving
the $c$-scaling hierarchy.
\end{quote}

\begin{thebibliography}{99}
\bibitem{paper18} \textit{Paper XVIII}, 2026.
\bibitem{paper19} \textit{Paper XIX}, 2026.
\bibitem{Kato1966} T.~Kato, \textit{Perturbation Theory for Linear Operators}, Springer, 1966.
\bibitem{Olver1997} F.~Olver, \textit{Asymptotics and Special Functions}, AK Peters, 1997.
\bibitem{Zettl2005} A.~Zettl, \textit{Sturm--Liouville Theory}, AMS, 2005.
\bibitem{ReedSimon2} M.~Reed, B.~Simon, \textit{Methods of Modern Mathematical Physics II}, Academic Press, 1975.
\bibitem{Osipov2013} A.~Osipov, V.~Rokhlin, H.~Xiao, \textit{Prolate Spheroidal Wave Functions}, Springer, 2013.
\bibitem{BGT1987} N.~H.\ Bingham, C.~M.\ Goldie, J.~L.\ Teugels, \textit{Regular Variation}, Cambridge UP, 1987.
\end{thebibliography}

\end{document}
